\documentclass[11pt]{article}

\usepackage{graphicx}
\usepackage{amsfonts}
\usepackage{amsmath,amssymb,mathtools}
\usepackage[latin1]{inputenc}
\usepackage{fancyhdr}
\usepackage{layout}
\usepackage{multicol}
\usepackage{enumitem}
\usepackage{xcolor}
\usepackage{framed}
\setcounter{MaxMatrixCols}{10}

\setlength{\voffset}{-1in}
\setlength{\topmargin}{1cm}
\setlength{\headheight}{1cm}
\setlength{\headsep}{1cm}
\setlength{\textheight}{24cm}
\setlength{\footskip}{1.2cm}
\setlength{\evensidemargin}{0in}
\setlength{\oddsidemargin}{0cm}
\setlength{\textwidth}{17cm}
\setlength{\marginparsep}{0cm}
\setlength{\marginparwidth}{0cm}
\pagestyle{fancy}
\rhead{Introduction to Group Theory - 1040158\\ Guangdong Technion}
\lhead{{\large\scshape Homework Assignment 1 \\ Winter - 2026}}
\cfoot{\thepage}
\renewcommand{\headrulewidth}{0.4pt}

\newcounter{ej}
\renewcommand{\theej}{\textbf{\arabic{ej}.}}
\newenvironment{ej}{
    \stepcounter{ej}
    \begin{list}{\theej}{
        \setlength{\itemindent}{0ex}
        \setlength{\labelwidth}{4ex}
        \setlength{\labelsep}{0.5ex}
        \setlength{\leftmargin}{4ex}
        \setlength{\listparindent}{0ex}
        \setlength{\topsep}{3mm plus 0.5mm minus 0.5mm}
        \setlength{\parsep}{0mm}
        \setlength{\itemsep}{1mm plus 0.3mm minus 0.3mm}
        \setlength{\rightmargin}{0ex}
    }
    \item
}{\end{list}}
\renewcommand{\labelenumi}{$($\alph{enumi}$)$ }

\definecolor{shadecolor}{gray}{0.92}
\newenvironment{sol}{
    \def\FrameCommand##1{\fcolorbox{black}{gray!20}{##1}}%
    \MakeFramed{\advance\hsize-\width \FrameRestore}%
    \noindent\textbf{Answer:} \ignorespaces
}{\endMakeFramed}

\newcommand{\Z}{\mathbb{Z}}
\newcommand{\Q}{\mathbb{Q}}

\begin{document}
    \thispagestyle{fancy}

    \begin{center}
        \large Autor: Rui Tang \\
        ID: 999025489
    \end{center}

    \begin{ej}
        Determine whether the following composition laws are associative,
        whether they are commutative and whether they have an identity:
        \begin{enumerate}
            \item $(x,y)\mapsto x-y$ on $\Z$;
            \item $(x,y)\mapsto x+y+xy$ on $\Z$; on $\Q\setminus\{-1\}$.
        \end{enumerate}
        Does any of them define a group structure? justify!

        \begin{sol}
            \textbf{(a)}

            \textbf{Associativity: No!}
            \begin{align*}
                ((x,y),z) &= (x-y,z)=(x-y)-z=x-y-z,\\
                (x,(y,z)) &= (x,y-z)=x-(y-z)=x-y+z.
            \end{align*}
            \[
                x-y-z\ne x-y+z\quad\text{in }\Z
            \]

            \textbf{Commutative: No!}
            \[
                (x,y)=x-y,\qquad (y,x)=y-x
            \]
            \[
                x-y\ne y-x\quad\text{in }\Z
            \]

            \textbf{Identity: No!}

            Suppose we have the identity called $e$. Then it satisfies, for
            every $x\in\Z$,
            \[
                (x,e)=x\quad\text{(1)},\qquad
                (e,x)=x\quad\text{(2)}.
            \]
            \begin{align*}
                \text{(1)}&\iff x-e=x\implies e=0,\\
                \text{(2)}&\iff e-x=x\implies x=0.
            \end{align*}
            It is contradiction to ``$\forall x\in\Z$''.

            So $\not\exists$ identity.

            \medskip
            \textbf{(b)}

            \textbf{Associativity: Yes! both on $\Z$ and $\Q\setminus\{-1\}$.}
            \begin{align*}
                ((x,y),z)
                    &= (x+y+xy,z)\\
                    &= (x+y+xy)+z+(x+y+xy)z\\
                    &= x+y+z+xy+yz+zx+xyz,\\[4pt]
                (x,(y,z))
                    &= (x,y+z+yz)\\
                    &= x+(y+z+yz)+x(y+z+yz)\\
                    &= x+y+z+xy+yz+zx+xyz.
            \end{align*}
            So $((x,y),z)=(x,(y,z))$.

            \textbf{Commutative: Yes! both on $\Z$ and $\Q\setminus\{-1\}$.}
            \[
                (x,y)=x+y+xy
            \]
            \[
                (y,x)=y+x+yx=x+y+yx
            \]
            \[
                \Downarrow\quad x,y,z\in\Z\text{ or }\Q\setminus\{-1\}
            \]
            So $(x,y)=(y,x)$.

            \textbf{Identity: the identity is $0$.}

            Since, for every $x\in\Z$ or $\Q\setminus\{-1\}$,
            \[
                \left\{
                \begin{aligned}
                    (x,0)&=x+0+x\cdot0=x,\\
                    (0,x)&=0+x+0\cdot x=x,
                \end{aligned}
                \right.
            \]
            \[
                \Rightarrow\ (x,0)=x=(0,x)
            \]
            So $0$ is the identity.

            \medskip
            Only (b) on $\Q\setminus\{-1\}$ determine a group structure.

            (a): can't because not associative.

            (b): $(x,y)\mapsto x+y+xy$ is $G\times G\longrightarrow G$.

            We only need inverse:
            \begin{align*}
                (x,x^{-1})=0
                    &\Rightarrow x+x^{-1}+xx^{-1}=0\\
                    &\Rightarrow x^{-1}=-\frac{x}{1+x}.
            \end{align*}
            Inverse exist only for $x\ne-1$.

            Therefore can't on $\Z$, only can on $\Q\setminus\{-1\}$.
        \end{sol}
    \end{ej}

    \begin{ej}
        Prove that a group $G$ that has the property that $a^2=e$ for all
        $a\in G$ is abelian.

        \begin{sol}
            proof:

            Since $\forall a\in G:\quad a^2=a\cdot a=e$,
            \[
                \text{then }a^{-1}(a\cdot a)=a^{-1}\cdot e
            \]
            \[
                \Rightarrow (a^{-1}\cdot a)a=a^{-1}
            \]
            \[
                \Rightarrow a=a^{-1}
            \]

            take $x,y\in G$
            \[
                xy=x^{-1}y^{-1}=(yx)^{-1}=yx
            \]
            So $G$ is abelian.
        \end{sol}
    \end{ej}

    \begin{ej}
        Let $a$ and $b$ be elements of a group $G$. Prove that the order of
        $ab$ is equal to the order of $ba$.

        \begin{sol}
            proof:

            Suppose the order of $ab$ is $n$.

            That is: $n$ is the minimal number s.t. $(ab)^n=e$.

            Now to prove the order of $ba$ is also $n$:

            \begin{itemize}
                \item To show $(ba)^n=e$:
                \begin{align*}
                    (ba)^n
                        &= (ba)(ba)\cdots(ba)\\
                        &= b(ab)(ab)\cdots(ab)a\\
                        &= b(ab)(ab)\cdots(ab)a(bb^{-1})\\
                        &= b(ab)(ab)\cdots(ab)(ab)b^{-1}\\
                        &= b(ab)^n b^{-1}\\
                        &= b\cdot e\cdot b^{-1}\\
                        &=e.
                \end{align*}

                \item To show $n$ is the smallest:

                Suppose $\exists m<n$ / $(ba)^m=e$.
                \begin{align*}
                    (ba)^m
                        &= (ba)(ba)\cdots(ba)\\
                        &= b(ab)(ab)\cdots(ab)a\\
                        &= b(ab)(ab)\cdots(ab)a(bb^{-1})\\
                        &= b(ab)(ab)\cdots(ab)(ab)b^{-1}\\
                        &= b(ab)^m b^{-1}.
                \end{align*}
                Then
                \begin{align*}
                    b(ab)^m b^{-1}&=e\\
                    \Rightarrow b(ab)^m b^{-1}b&=e\cdot b\\
                    \Rightarrow b(ab)^m&=b\\
                    \Rightarrow b^{-1}b(ab)^m&=b^{-1}b\\
                    \Rightarrow (ab)^m&=e.
                \end{align*}
                It is contradiction to order $(ab)=n$ $(m<n)$.
            \end{itemize}

            Therefore: order $(ba)=n$.
        \end{sol}
    \end{ej}

    \begin{ej}
        Let $G$ be a group. Prove that the following subsets of $G$ are
        subgroups.
        \begin{enumerate}
            \item $H_1\cap H_2$ where $H_i\le G$.
            \item The \emph{centraliser} of $a\in G$ defined by
                $C_G(a)=\{x\in G\mid ax=xa\}$.
            \item The \emph{centre} of $G$ defined by
                $Z(G)=\{x\in G\mid xy=yx,\ \forall y\in G\}$.
        \end{enumerate}

        \begin{sol}
            \textbf{(a)}

            proof: we have the property that:

            \textcolor{red}{Property: $H\le G$ iff}
            \[
                \textcolor{red}{\text{\textcircled{1}}\ e\in H,\qquad
                \text{\textcircled{2}}\ \forall a,b\in H\Rightarrow ab\in H,
                \qquad
                \text{\textcircled{3}}\ \forall a\in H\Rightarrow a^{-1}\in H.}
            \]
            So we need to show \textcircled{1}\textcircled{2}\textcircled{3}
            with $H_1\cap H_2$.

            \textcircled{1}: since $H_1\le G,\ H_2\le G$,
            \[
                \text{then }e\in H_1,\ e\in H_2
                \quad\Rightarrow\quad e\in H_1\cap H_2.
            \]

            \textcircled{2}: take $a,b\in H_1\cap H_2$.
            \[
                \left\{
                \begin{aligned}
                    a,b\in H_1&\xRightarrow{H_1\le G}ab\in H_1,\\
                    a,b\in H_2&\xRightarrow{H_2\le G}ab\in H_2.
                \end{aligned}
                \right.
            \]
            Then $ab\in H_1\cap H_2$.

            \textcircled{3}: take $x\in H_1\cap H_2$.

            In particular $x\in H_1\ \wedge\ x\in H_2$.
            \[
                \Rightarrow x^{-1}\in H_1,\ x^{-1}\in H_2
                \quad\Rightarrow\quad x^{-1}\in H_1\cap H_2.
            \]
            Therefore $H_1\cap H_2\le G$.

            \medskip
            \textbf{(b)}

            proof: the same, we need to show \textcircled{1}\textcircled{2}
            \textcircled{3} with $C_G(a)$.

            \textcircled{1}: since $e\cdot a=a=a\cdot e$

            then $e\in C_G(a)$.

            \textcircled{2}: take $x,y\in C_G(a)$.
            Then
            \[
                \left\{
                \begin{aligned}
                    xa&=ax\quad\text{(1)},\\
                    ya&=ay\quad\text{(2)}.
                \end{aligned}
                \right.
            \]
            \begin{align*}
                (xy)a&=x(ya)=x(ay)\\
                    &=(xa)y\\
                    &=(ax)y\\
                    &=a(xy).
            \end{align*}
            So $(xy)\in C_G(a)$.

            \textcircled{3}: take $x\in C_G(a)$.
            \[
                \text{then }xa=ax\quad\text{(1)}
            \]
            \[
                \text{(1)}\Rightarrow a=x^{-1}ax
            \]
            \begin{align*}
                \text{then }ax^{-1}
                    &=(x^{-1}ax)x^{-1}\\
                    &=(x^{-1}a)(xx^{-1})\\
                    &=x^{-1}a.
            \end{align*}
            So $x^{-1}\in C_G(a)$.

            Therefore $C_G(a)\le G$.

            \medskip
            \textbf{(c)}

            \textcircled{1}: Since $e\cdot y=y=y\cdot e\quad\forall y\in G$

            then $e\in Z(G)$.

            \textcircled{2}: take $a,b\in Z(G)$.

            Then
            \[
                \left\{
                \begin{aligned}
                    ax&=xa\quad\forall x\in G\quad\text{(1)},\\
                    by&=yb\quad\forall y\in G\quad\text{(2)}.
                \end{aligned}
                \right.
            \]
            \begin{align*}
                (ab)z&=a(bz)=a(zb)=(az)b\\
                    &\overset{\text{(1)}}{=}(za)b\\
                    &=z(ab).
            \end{align*}
            Then $ab\in Z(G)$.

            \textcircled{3}: take $a\in Z(G)$.

            Then $ax=xa\quad\forall x\in G$.
            \begin{align*}
                a^{-1}x
                    &=a^{-1}(axa^{-1})\\
                    &=(a^{-1}a)(xa^{-1})\\
                    &=xa^{-1}\quad\forall x\in G.
            \end{align*}
            So $a^{-1}\in Z(G)$.

            Therefore $Z(G)\le G$.
        \end{sol}
    \end{ej}

    \begin{ej}
        Prove that the inverse of a permutation matrix is its transpose.

        \begin{sol}
            \textcolor{orange}{Def: permutation matrix is a $n\times n$ matrix
            that every row, column has only one ``1''.}

            proof: Suppose we have a permutation matrix $P$.

            To show $P^{-1}=P^T$ we just need to show that
            $PP^T=I=P^TP$.

            \textcircled{1}\quad $PP^T=I$:

            Suppose $P\in M_{n\times n}$,
            \[
                P=(P_{ij})=
                \begin{bmatrix}
                    P_{11}&P_{12}&\cdots&P_{1n}\\
                    P_{21}&P_{22}&\cdots&P_{2n}\\
                    \vdots&\vdots&\ddots&\vdots\\
                    P_{n1}&P_{n2}&\cdots&P_{nn}
                \end{bmatrix},
                \qquad
                P^T=
                \begin{bmatrix}
                    P_{11}&P_{21}&\cdots&P_{n1}\\
                    P_{12}&P_{22}&\cdots&P_{n2}\\
                    \vdots&\vdots&\ddots&\vdots\\
                    P_{1n}&P_{2n}&\cdots&P_{nn}
                \end{bmatrix}.
            \]
            Then
            \[
                (PP^T)_{ij}=P_{i1}P_{j1}+P_{i2}P_{j2}+\cdots+P_{in}P_{jn}.
            \]

            If $i\ne j$:
            \[
                (PP^T)_{ij}
                =P_{i1}P_{j1}+P_{i2}P_{j2}+\cdots+P_{in}P_{jn}=0.
            \]

            If $i=j$:
            \begin{align*}
                (PP^T)_{ij}
                    &=P_{i1}P_{i1}+P_{i2}P_{i2}+\cdots+P_{in}P_{in}\\
                    &=P_{i1}^2+P_{i2}^2+\cdots+P_{in}^2
                    \quad\textcolor{green}{(One of them = 1, other = 0)}\\
                    &=1.
            \end{align*}
            So $PP^T=I$.

            Similarly: $P^TP=I$.

            Therefore: $P^{-1}=P^T$.
        \end{sol}
    \end{ej}

\end{document}